% THEORY.tex — payout mechanisms as operators on attention distributions.
% Companion to THEORY.md (Russian source of truth); English for journal use.
% v0.4: editorial pass — math untouched, two-layer typography (route
% paragraphs in italics; proofs in small type). Compile: tectonic THEORY.tex
\documentclass[11pt]{article}
\usepackage[T1]{fontenc}\usepackage[utf8]{inputenc}
\usepackage{amsmath,amssymb,amsthm,geometry,booktabs,etoolbox}
\geometry{margin=1.05in}
% ---- two-layer typography: proofs small, route paragraphs italic ----
\AtBeginEnvironment{proof}{\small}
\newcommand{\why}[1]{\par\medskip\noindent{\itshape\textbf{Why this matters.} #1}\par\medskip}
\newcommand{\whatitsays}[1]{\par\medskip\noindent{\itshape\textbf{What it says.} #1}\par\medskip}
% numbering matches the THEORY.md canon (T1, L2, T4, corollaries (a)-(d))
\newtheorem*{thmone}{Theorem 1}
\newtheorem*{thmfour}{Theorem 4}
\newtheorem*{lemtwo}{Lemma 2}
\newtheorem*{cora}{Corollary (a)}
\newtheorem*{corb}{Corollary (b)}
\newtheorem*{corc}{Corollary (c)}
\newtheorem*{cord}{Corollary (d)}
\newtheorem{assumption}{Assumption}
\theoremstyle{remark}\newtheorem*{remark}{Remark}
\newcommand{\E}{\mathbb{E}}
\newcommand{\Pbar}{\bar{P}}
\title{Payout Mechanisms as Operators on Attention Distributions:\\
Identities, Invariances, and Tail Action\vspace{-2pt}}
\author{Tonify Research}
\date{August 2026 $\cdot$ v0.4 (editorial pass over v0.3; math untouched) $\cdot$ companion code: \texttt{tonify-sims}}
\begin{document}\maketitle

\paragraph{How to read this paper.}
A streaming platform pools subscriber money and splits it among artists;
the three ways to split are \emph{pro-rata} (one pool, paid by global play
share), \emph{user-centric} (each wallet split only across its owner's own
plays), and \emph{direct} (no pool: fans pay artists directly). This paper
treats all three as operators mapping a play matrix to artist incomes and
answers four questions: who gains from a pro-rata\,$\to$\,user-centric
switch (Theorem~1); whether any pool-rule reform can touch the
signed-vs-independent gap (Lemma~2); whether the direct economy reduces
top-end concentration (it does not --- it relocates inequality to a mass of
exact zeros); and whether direct income can ever stochastically dominate
the pool (it cannot --- Theorem~4). Every claim below is machine-verified
in the companion code and survived two adversarial red-team passes. The
italicized paragraphs (``Why this matters'' / ``What it says'') form a
complete non-technical route: a reader may skip every formula and proof
(set in small type) and still carry away the results.

\section{Setup and the generative model of sim4}
Artists $i\in N$, listeners $u\in U$ ($|U|=U$), integer play matrix
$p=(p_{iu})\in\mathbb{Z}_+^{N\times U}$ over one period. Row/column sums:
$P_i=\sum_u p_{iu}$ (the artist's plays), $P_u=\sum_i p_{iu}$ (the
listener's intensity), $T=\sum_u P_u$, mean intensity
$\Pbar=T/U$. Equal wallets $W>0$; pool $R=UW$.

\begin{assumption}[A1, load-bearing]\label{a1}
Every paying subscriber listens: $P_u\ge 1$. Silent payers are excluded from
\emph{both} pools; otherwise budget balance below fails by exactly
$W\cdot U_{\mathrm{sil}}$ and all ratios acquire the factor
$U_{\mathrm{act}}/U_{\mathrm{tot}}$.
\end{assumption}

\paragraph{Mechanisms.} Pro-rata $\mathrm{PR}_i = UW\,P_i/T$; user-centric
$\mathrm{UC}_i = W\sum_u p_{iu}/P_u$; direct
$D_i=\tau\sum_{f=1}^{S_i}\sum_{t=1}^{k}G_{ft}$ with
$S_i\,|\,A_i\sim\mathrm{Bin}(A_i,\sigma)$, audience $A_i$ (the number of
listeners with $p_{iu}>0$), superfan share $\sigma$, i.i.d.\ gifts
$G\ge 0$, $\E G=\bar g$, all moments finite; artist share $\tau$. A label
contract acts multiplicatively after the rule: $\rho\cdot F_i(p)$.

\paragraph{Generative model (sim4).} Artist popularity: piecewise
$s_i\sim$ (lognormal body $/$ log-bridge $/$ Pareto($\alpha{=}1.4$) tail),
calibrated to anchors $T_1{:}\,87\%$ of artists below $10^3$ streams/yr,
$T_2{:}\,2.6\%$ above $225{,}734$, $T_3{:}$ top $0.28\%$ holding
$40$--$55\%$ of streams (gate band; the seed-42 matrix realizes $40.3\%$,
sim1's larger world $44.5\%$); target audiences $\ell_i=s_i/\mathrm{PL}$,
$\mathrm{PL}=21.21$ plays per pair (LFM-1b). Playlist sizes
$K_u\sim\mathrm{LogNormal}(\ln 12,\,1.0)$ (assumed; an MVP-measured dial).
Artist selection: exact weighted sampling without replacement,
$\mathbb{P}(i\in S_u)\propto \ell_i$ (Gumbel top-$K_u$), i.e.\ preferential
attachment. Pair intensity, three named regimes on one graph:
\[
\text{(a) ergodic: integer } p_{iu},\ P_u\equiv 2048;\qquad
\text{(b) } p_{iu}=\lceil\mathrm{LN}(\ln 5.16,\,1.676)\rceil
\text{ rescaled to } s_i;\qquad
\text{(c) } p_{iu}\propto \mathrm{LN}\cdot(A_i/\bar A_g)^{\gamma}.
\]
Regime (a) has $P_u\equiv 2048$ exactly (a power of two: every
$p_{iu}/P_u$ is a dyadic rational, so the control identity is exact in
floating point) and is the control (Corollary (a));
$\gamma$ dials the coupling between pair intensity and audience size ---
its role becomes central in the tail-action analysis.

\section{The ratio identity}

\why{Debates about user-centric run on intuition (``it helps niche
artists''). The identity below replaces intuition with a formula: an
artist's user-centric-to-pro-rata income ratio is a specific average over
that artist's own audience, so the winners of a reform can be computed
before the reform.}

\begin{thmone}[Identity]
For any artist $i$ with $P_i>0$, under equal wallets,
\[
\frac{\mathrm{UC}_i}{\mathrm{PR}_i}
=\E_{w^i}\!\left[\frac{\Pbar}{P_u}\right],
\qquad w^i_u=\frac{p_{iu}}{P_i}.
\]
\end{thmone}
\begin{proof}
$\mathrm{UC}_i/\mathrm{PR}_i
=\bigl(W\sum_u p_{iu}/P_u\bigr)\big/\bigl(UW P_i/T\bigr)
=(T/U)\sum_u (p_{iu}/P_i)(1/P_u)=\Pbar\,\E_{w^i}[1/P_u]$.
\end{proof}

\noindent With $m_i=\E_{w^i}[P_u]$ (arithmetic) and
$H_i=1/\E_{w^i}[1/P_u]$ (harmonic mean of listener intensity under play
weights),
\[
\frac{\mathrm{UC}_i}{\mathrm{PR}_i}=\frac{\Pbar}{H_i}
=\underbrace{\frac{\Pbar}{m_i}}_{\text{level}}\cdot
\underbrace{\frac{m_i}{H_i}}_{\text{dispersion}\ \ge 1},
\]
the second factor $\ge 1$ by AM--HM, with equality iff $P_u$ is constant on
the support of $w^i$.

\whatitsays{The ratio splits into level $\times$ dispersion: a ``light''
audience (below-average listeners) helps through the first factor, and a
\emph{heterogeneous} audience always helps through the second. So the folk
claim ``user-centric helps the niche'' is wrong in an instructive way: what
helps is not nicheness but lightness and heterogeneity of one's listeners.}

\begin{cora}[Coincidence; the naive converse is false]
$P_u\equiv\Pbar$ implies $\mathrm{UC}=\mathrm{PR}$ componentwise. The
converse fails: with one artist and two listeners, $p=(10,20)$,
intensities differ yet $\mathrm{UC}_1=\mathrm{PR}_1=2W$. Exact criterion:
$\mathrm{UC}_i=\mathrm{PR}_i$ for all $i$ iff the vector
$c_u=1/P_u-U/T$ lies in $\ker p$.
\end{cora}

\begin{corb}[Zero sum]
$\sum_i \mathrm{UC}_i=\sum_i\mathrm{PR}_i=UW$: user-centric is a pure
redistribution of the same pool.
\end{corb}

\begin{corc}[Crossover]
$\mathrm{UC}_i>\mathrm{PR}_i$ iff $H_i<\Pbar$. On type populations this
reduces to the condition of Alaei, Makhdoumi, Malekian and Peke\v c
(\emph{Mgmt.\ Sci.}\ 2022, Props.\ 2, 8): artist $j$ weakly prefers
pro-rata iff $\sum_i q_i\pi_{ij}(\lambda_i-\bar\lambda)\ge 0$.
\end{corc}

\begin{cord}[Dispersion premium]
Fix $m_i=\Pbar$. Then $\mathrm{UC}_i/\mathrm{PR}_i=m_i/H_i\ge 1$, with
equality iff audience intensity is degenerate: at the same mean level,
\emph{within-audience dispersion strictly raises} user-centric income.
\end{cord}

\section{Pass-through invariance}

\why{The central policy dispute: is the artist squeezed by the pool rule or
by the label contract? If a signed artist keeps seven cents on the dollar,
can any reform of the splitting rule close that gap? The lemma answers:
never in income terms --- and sometimes in minimum-viable-audience terms.}

\begin{lemtwo}[Two-part invariance]
\emph{(2a)} For any pool operator $F$ \emph{applied at the artist level},
with the multiplicative contract strictly downstream and on the domain
$F_i(p)>0$: $\rho F_i/F_i=\rho$ --- the signed/independent \emph{income} gap
is invariant to the rule (a nonlinear $F$ applied to a licensor's
\emph{catalogue} breaks this). \emph{(2b)} The
\emph{minimum-viable-audience} gap satisfies
$\mathrm{MVA}_s/\mathrm{MVA}_\iota=1/\rho$ iff
$F^{-1}(V/\rho)=F^{-1}(V)/\rho$; linearity of $F$-income in audience is
sufficient but \emph{not} necessary (a strictly increasing log-periodic
$F$ is nonlinear yet yields $1/\rho$ for all $V$). A Deezer-style
$\times 2$ boost above a threshold yields $1/(2\rho)$ instead
(counterexample in companion code) while the income ratio stays $\rho$
pointwise.
\end{lemtwo}

\whatitsays{A rule reform cannot move the contractual income gap by one cent;
it can move the ``minimum audience to make a living'' gap, but only by
breaking the scaling criterion above. This formalizes the empirical
regularity that label shares did not move across rule changes (Pedersen
2020). Breaking channels: endogenous $\rho(F)$; additive components; payout
thresholds (accumulation delays $\times 1/\rho$ in time); behavioral
endogeneity of $p$; component/market-heterogeneous
$\rho_{\mathrm{eff}}=\sum_m\rho_m w_{im}(F)$; payee-level aggregation of a
nonlinear rule.}

\section{Tail action and the atom at zero}\label{tail}

\why{Streaming's chief indictment is concentration: the top percent takes
nearly everything. Would a direct fan-to-artist economy spread income more
evenly? This section answers for the top and the bottom of the
distribution separately --- and the answer disappoints every mechanism.}

Let $\mathbb{P}(A>x)=x^{-\alpha}L(x)$. PR is linear: index preserved.
UC: $\mathrm{UC}_i=WP_i\,\E_{w^i}[1/P_u]$ with bounded multiplier;
under tail-conditional concentration of the multiplier the index is
preserved, and heavier tails are impossible for any dependence.
Direct: binomial thinning plus a light-tailed compound sum preserve
$\alpha$ (Jessen--Mikosch 2006, Lem.\ 3.7; Fa\"y et al.\ 2006;
Denisov--Foss--Korshunov 2010; Breiman 1965 enters only through the
multiplicative contract layer), while
$\mathbb{P}(D_i=0\,|\,A_i)\ge(1-\sigma)^{A_i}\ge e^{-1}(1-\sigma)$ for
$A_i\le 1/\sigma$: an \emph{atom at zero} on small audiences. With zero
mass $q$, $G_{\mathrm{direct}}=q+(1-q)G^{+}$: the mechanism relocates
inequality from the intensive to the extensive margin; no mechanism
softens the upper tail.

\whatitsays{No mechanism touches how much the top takes; what changes is the
\emph{shape of the bottom}. Pool rules pay many tiny positive amounts;
the direct economy converts most of those into exact zeros --- at the
calibrated superfan share, at least a third of small-audience artists earn
literally nothing. Inequality is relocated, not reduced.}

\section{Dominance blockade}

\why{Suppose the direct economy pays four times the pool on average. Should
a small artist leave the pool? Not necessarily: the average hides a large
chance of zero. This section compares the two incomes by their whole
distributions, in the standard language of stochastic dominance.}

Fix $A$ with $(1-\sigma)^A>0$ and pooled income $X_A=cA>0$
(deterministic given $A$); $Y_A=D_A$ as above.

\begin{thmfour}[Full blockade]
(i) $Y_A$ never first-order dominates $X_A$;
(ii) $Y_A$ never second-order dominates $X_A$, for any parameters;
(iii) $X_A\succeq_2 Y_A$ iff $\E Y_A\le X_A$, i.e.\ iff
$k\le d^\ast=\mathrm{PL}\cdot\mathrm{rate}/(\sigma\bar g\tau)$;
(iv) for $k>d^\ast$ the pair is incomparable in both orders.
\end{thmfour}
\begin{proof}
(i) For $x\in(0,X_A)$: $\mathbb{P}(X_A>x)=1$ while
$\mathbb{P}(Y_A>x)\le 1-(1-\sigma)^A<1$.
(ii) SSD against degenerate $X_A$ requires
$\int_0^t F_Y\le\max(0,t-X_A)$; at $t=X_A$ the left side is
$\ge(1-\sigma)^A X_A>0$.
(iii) ($\Leftarrow$) $\int_0^t F_Y=t-\E\min(Y,t)\ge t-\E Y\ge t-X_A$ and
$\ge 0$. ($\Rightarrow$) $t\to\infty$ gives $\E Y\le X_A$.
(iv) From (i)--(iii).
\end{proof}
\begin{remark}
The correct comparison is mean versus risk: for $k>d^\ast$ direct pays an
expectation premium priced by small-audience lottery risk;
$\mathbb{P}(Y_A<X_A)\le\inf_{t>0}e^{tX_A}\bigl(1-\sigma+\sigma\varphi(t)^k
\bigr)^A$ with $\varphi(t)=\E e^{-t\tau G}<1$ --- exponential decay in $A$.
In the companion sim1 calibration at $A{=}30$, $k{=}4$: mean advantage
$\times 4$, yet $\mathbb{P}(Y<X)\approx 0.60$ --- and on the full 2M-sample
run $\mathbb{P}(Y<X)-\mathbb{P}(Y=0)=0$ to machine precision: losing to
the pool is exactly the zero outcome.
\end{remark}

\whatitsays{The zero atom blocks stochastic dominance in both directions and
in both orders, for all parameter values: the direct economy can never be
unconditionally better than the pool, no matter how generous the fans. The
honest frame is mean versus risk --- direct pays an expectation premium
priced by a lottery on small audiences, and only raising the superfan
conversion $\sigma$ (not payment frequency $k$) attacks the lottery
itself.}

\end{document}
